% ============================================================
% РЕФЕРАТ
% ============================================================
\structure{РЕФЕРАТ}
\addcontentsline{toc}{structure}{Реферат}

\textbf{Объект исследования} --- инфраструктура бессерверного инференса моделей машинного обучения, развернутая в среде Kubernetes с применением механизма масштабирования до нуля (scale-to-zero).

\textbf{Предмет исследования} --- методы кэширования, доставки и виртуализации доступа к файлам весов ML-моделей при холодном старте инференс-подов в Kubernetes-кластере.

\textbf{Цель работы} --- исследование и разработка архитектуры распределённого кэширования файлов моделей машинного обучения для бессерверного инференса в Kubernetes, позволяющей сократить задержку холодного старта, уменьшить нагрузку на внешние источники данных и повысить доступность системы при горизонтальном масштабировании.

Для достижения поставленной цели решены следующие задачи: проведён анализ причин холодного старта и влияния роста размеров ML-моделей на инфраструктуру инференса; исследованы существующие распределённые решения для кэширования моделей --- Alluxio, JuiceFS, Dragonfly, CubeFS, Rook/Ceph; разработан и обоснован промежуточный подход на базе NFS с механизмами координации параллельных загрузок через файловые блокировки и Kubernetes Lease; спроектирована и частично реализована специализированная система распределённого chunk-кэша, включающая агент \texttt{storage-agent}, работающий как DaemonSet на каждой worker-ноде, FUSE-клиент \texttt{storage-mounter}, предоставляющий ML-runtime стандартный файловый интерфейс, и Redis-кластер в роли реестра метаданных чанков и активных агентов; сформирована методика оценки эффективности предложенной архитектуры.

\textbf{Ключевые методы}: трёхуровневое кэширование по схеме «локальный диск -- соседняя нода -- внешний репозиторий»; FUSE-монтирование виртуальной файловой системы; P2P-доставка чанков между агентами по TCP; локальная связь через Unix Domain Socket; Redis heartbeat как реестр метаданных; атомарная запись чанков через временный файл и \texttt{os.Rename}.

\textbf{Ожидаемые результаты}: снижение количества повторных загрузок весов из внешнего репозитория, уменьшение нагрузки на HuggingFace~Hub, ускорение холодного старта при наличии локального или соседнего кэша, устранение NFS как единственной точки отказа.

\textbf{Ключевые слова}: холодный старт, машинное обучение, бессерверный инференс, Kubernetes, кэширование, P2P, FUSE, Redis, chunk, storage-agent, storage-mounter.
